% SPDX-FileCopyrightText: 2026 Will Cook
% SPDX-License-Identifier: CC-BY-4.0
% System-architecture paper; build with tectonic.
\documentclass[10pt]{article}
\usepackage{paper-house-style}
\usepackage{booktabs,array,float,placeins,multicol}
\usepackage{tikz}
\usetikzlibrary{arrows.meta,positioning,calc,fit,shapes.geometric}
\usepackage[colorlinks=true,linkcolor=linkinternal,urlcolor=linkexternal,citecolor=linkinternal]{hyperref}
\hypersetup{
  bookmarksnumbered=true,bookmarksopen=true,bookmarksopenlevel=1,
  pdftitle={Problem-Sized Lean Worlds},
  pdfauthor={Will Cook},
  pdfsubject={Architecture and trust boundaries of an AI-assisted Lean research system},
  pdflang={en-GB},
  pdfkeywords={formal mathematics, Lean 4, AI agents, research architecture, claim verification, reproducibility},
}
\newcommand{\repobase}{https://github.com/wcook04/plectis-lean-erdos249-257}
% BEGIN generated_semantic_coverage_macros
\newcommand{\SemanticAuthoredTheoremLike}{144,733}
\newcommand{\SemanticAuthoredInterpreted}{139,817}
\newcommand{\SemanticDirectEvidence}{3,329}
\newcommand{\SemanticContextual}{136,488}
\newcommand{\SemanticStructuralOnly}{4,916}
\newcommand{\SemanticAuthoredPercent}{96.6}
% END generated_semantic_coverage_macros
\newcommand{\repolink}[2]{\href{\repobase/blob/ca0e13f8acf5ccf48506e4bdb870953d3a0856fa/#1}{\texttt{#2}}}
\emergencystretch=1.7em

\runninghead{Problem-sized Lean worlds}
\title{Problem-Sized Lean Worlds}
\subtitle{An authority-separated architecture from AI search to public mathematical claims}
\author{Will Cook}
\date{August 2026}

\begin{document}
\maketitle
\paperprovenance
\fontsize{9.2pt}{11pt}\selectfont

\begin{abstract}
This paper presents an executable architecture for AI-assisted mathematics
under proof abundance. It separates six things that are commonly collapsed:
cognition authority, permission to change shared state, validation capacity,
evidence class, public-claim authority, and reviewer attention. Work advances
only when the next gate matches the evidence. A failed route may change later
search or validation without being relabelled as a theorem.

The architecture organises work around \emph{problem-sized mathematical
worlds}, not isolated theorem files. Each world joins formal declarations,
experiments, source literature, failed mechanisms, open obligations, and a
reviewed public boundary. Formal dependencies, authored mathematical meaning,
and public claims remain separate graphs, so a dense connection cannot make a
claim stronger. An agent enters a world through a compiled comprehension
packet: it federates the corpora by content digest, states the endpoint and
claim ceiling before any route, separates proved results from open producers
and closed routes, and declares both what it omitted and when it exceeded its
context budget. The contribution is an implemented architecture that connects
research state, scoped changes, formal checks, and public explanation while
keeping their evidential roles distinct; the related-work section compares it
with neighbouring systems along stated dimensions and claims no priority.

The mathematical workflow keeps conjecture, computation, formal proof,
interpretation, public wording, and external review distinct. Computation may
reject a route or supply finite evidence; it does not prove an unbounded
theorem. Lean verifies that a proof establishes the formal statement written in
the source; it does not verify whether that statement captures the intended
mathematics or whether the paper describes it
well. Comparator checks selected propositions and their axiom boundary; a
local review-selection layer prepares result families for external
registration and review. Neither decides novelty, importance, acceptance, or
whether an open problem has been solved.

The public repository is a self-contained output whose Lean source, claim
record, papers, navigation, and release checks can be inspected from a fresh
clone. A worked example follows a finite certificate for Erd\H{o}s Problem 249
and shows why a checked finite range cannot become an unbounded conclusion.
This is a working prototype rather than a validated multi-user service. As of
31 August 2026, the author had not recorded a completed external cold-clone use
or an accepted external contribution.
The rule is simple: each stage may pass forward only the claim its evidence
supports, together with the boundary it does not cross. All eight problems
remain open.
The paper does not claim a solution to any of them.
\end{abstract}

\begin{paperresult}{Main contribution and claim ceiling}
\textbf{Contribution.} The prototype keeps six authorities separate and
makes every transition from conjecture to public wording carry its evidence
and its limit.  Problem-sized worlds join formal dependencies, mathematical
interpretation, failed routes, and reviewed claims without treating any one
graph as authority for the others.  \textbf{Demonstration.} The mechanism is
exercised on a self-contained public Lean corpus and an audited finite
certificate for Erd\H{o}s Problem 249.  \textbf{Ceiling.} No open problem is
claimed solved, and external use remains untested.
\end{paperresult}

\section{The problem: many kinds of evidence}

Suppose an AI system proposes a proof of a mathematical statement. Several
different questions immediately arise.

\begin{enumerate}
\item Did the system understand the intended problem?
\item Did it explore the important alternatives and notice counterexamples?
\item Does the proposed formal statement say what the informal statement says?
\item Does Lean accept a proof of that exact formal statement?
\item Does the public explanation stay within the proved scope?
\item Has anyone outside the producing system reviewed, accepted, or absorbed
      the result?
\end{enumerate}

No single test answers all six. A numerical experiment can expose a false
conjecture without proving a universal one. Lean can certify a theorem while
remaining silent about the prose around it. A release checker can preserve a
reviewed sentence without making the review correct. External attention can
signal interest without validating a proof. The architecture exists to keep
these questions separate while allowing work to move between them.

The \href{\repobase}{public repository studied here} contains Lean source,
papers, and release machinery around eight open Erd\H{o}s problems. At this
revision all eight problems remain open. The repository contains substantial
intermediate theorems, exact reformulations, conditional reductions, finite
certificates, and no-go results. The surrounding private workbench is broader:
it supports research, agent coordination, computational experiments,
formalisation, exposition, and controlled public projection. Private machinery
explains how the work was produced; it supplies no hidden proof authority to
the public clone.

\paragraph{The novelty claim and its ceiling.}
This paper does not claim to have invented agents, queues, file locks, theorem
graphs, Lean checking, pull requests, or credit records. Its candidate
contribution is an executable \emph{claim-transition architecture}. It models
reasoning, mutation, validation, interpretation, publication, and review as
separately scarce and separately authorised operations. The operations are
recombined only at explicit fan-in gates. The same machinery carries proofs,
counterexamples, no-gos, and unresolved obligations while preserving their
different evidence classes. Crucially, a local failure may alter the next
agent's route, context, lease, experiment, or validator, but never the truth
status of a mathematical statement.

Four status words are used deliberately. \emph{Implemented} means source and
an executable check or receipt exist. \emph{Projection} means a generated view
over more authoritative records. \emph{Inactive} means implemented machinery
was not running at the reported snapshot. \emph{Hypothesis} means a proposed
experiment, not a reported capability. Thus the router, work leases, corpus
maps, Lean gates, Comparator, review-selection records, and release checks are
implemented; dashboards and graph views are projections; the resident
maintenance daemon was inactive at one recorded snapshot; and mass
frontier-model mining or training on the no-go graph remain hypotheses.

\subsection{Two architectural claims}

\paragraph{A problem is a mathematical world, not a folder.}
For an agent, the useful unit is neither the repository nor one theorem. It is
a bounded neighbourhood inside a problem-sized world. The neighbourhood begins
with the exact endpoint and current claim ceiling, then selects established
premises, nearby declarations, open producers, consumers, alternative
coordinates, executable experiments, scoped falsifiers and no-gos, source
literature, and public boundaries. Every returned edge says why it exists and
what authority it lacks; the packet says what it omitted and where to expand.
This makes deep corpora navigable without converting file proximity, lexical
similarity, or model interpretation into proof.

\paragraph{Authority is separated before it is recombined.}
The architecture treats six resources as non-fungible. More reasoning cannot
buy a write lease; a write cannot buy a Lean receipt; a Lean receipt cannot buy
semantic correctness; approved wording cannot buy novelty or community
acceptance. They meet only through typed artefacts at an explicit gate.

\begin{table}[!ht]
\centering\scriptsize
\begin{tabular}{@{}>{\raggedright\arraybackslash}p{0.19\linewidth}
                    >{\raggedright\arraybackslash}p{0.30\linewidth}
                    >{\raggedright\arraybackslash}p{0.39\linewidth}@{}}
\toprule
\papertablehead
Resource & Governing object & What possession cannot imply \\
\midrule
Reasoning scope & Task-conditioned context and trace & Mutation permission or mathematical truth \\
Mutation permission & Exact path/work lease and change & Formal acceptance or public promotion \\
Validator capacity & Single-flight slot and build receipt & Theorem failure when a request is deferred \\
Evidence class & Experiment, counterexample, theorem, or authored relation & Automatic promotion into another class \\
Public-claim permission & Reviewed claim record and release relationship & Novelty, independent review, or acceptance \\
Reviewer attention & Ranked packet and recorded human outcome & Proof authority or canonical status \\
\bottomrule
\end{tabular}
\caption{The six non-fungible resources of the claim-transition architecture.}
\label{tab:nonfungible}
\end{table}

There are consequently two coupled graphs with a guarded crossing. Mathematical
evidence changes the mathematical graph: theorem, counterexample, no-go, or
bounded experiment, each with its own class. Operational evidence changes the
control graph: a route miss, stale view, repeated workaround, validation
failure, or resource conflict may justify a new router, skill, check, or
standard after a generalisation guard. Reviewed claim and publication artefacts
connect the two. Neither graph may rewrite the other by implication.

\section{The whole lifecycle in one picture}\label{sec:lifecycle}
\papersectiontarget{systems-lifecycle}

Figure~\ref{fig:lifecycle} shows the main path. It is a loop rather than a
one-way publishing pipeline: failed proofs, reviewer objections, and public
drift can all return a precise lesson to an earlier stage. What may move
backward is information about an error or a better method. What may not move
backward is authority: a paper cannot make an unproved statement true, and an
external review route cannot alter what Lean checked.

\begin{figure}[!t]
\centering
\begin{tikzpicture}[
  box/.style={paper stage, text width=2.05cm},
  private/.style={box, fill=MidnightBlue!5, draw=MidnightBlue!55},
  formal/.style={box, fill=TealBlue!6, draw=TealBlue!60!black},
  public/.style={box, fill=ForestGreen!6, draw=ForestGreen!55!black},
  assure/.style={box, fill=BurntOrange!7, draw=BurntOrange!65!black},
  arr/.style={paper arrow},
  back/.style={paper feedback},
  lab/.style={font=\scriptsize\sffamily\bfseries, text=gray!70!black}
]
\node[private] (intent) at (0,0) {{\scriptsize\sffamily\bfseries 1 \; SELECT}\\[3pt]
  \textbf{Intent and route}\\[2pt]\scriptsize question, scope, claim};
\node[private, right=3.5mm of intent] (reason) {{\scriptsize\sffamily\bfseries 2 \; TEST}\\[3pt]
  \textbf{Reason and test}\\[2pt]\scriptsize candidates and failures};
\node[formal, right=3.5mm of reason] (lean) {{\scriptsize\sffamily\bfseries 3 \; CHECK}\\[3pt]
  \textbf{Formalise}\\[2pt]\scriptsize proof and kernel verdict};
\node[public, right=3.5mm of lean] (claim) {{\scriptsize\sffamily\bfseries 4 \; EXPLAIN}\\[3pt]
  \textbf{Public claim}\\[2pt]\scriptsize result, evidence, limit};
\node[assure, right=3.5mm of claim] (review) {{\scriptsize\sffamily\bfseries 5 \; RELEASE}\\[3pt]
  \textbf{Assure}\\[2pt]\scriptsize replay and boundary checks};
\draw[arr] (intent) -- (reason);
\draw[arr] (reason) -- (lean);
\draw[arr] (lean) -- (claim);
\draw[arr] (claim) -- (review);
\draw[back] (review.south) to[out=-115,in=-65,looseness=1.15]
  node[lab, below=4pt] {failures and review return as evidence; authority stays put} (reason.south);
\node[font=\scriptsize\sffamily\bfseries, text=MidnightBlue!75!black,
  above=5mm of $(intent.north)!0.5!(reason.north)$] {PRIVATE WORKBENCH};
\node[font=\scriptsize\sffamily\bfseries, text=ForestGreen!62!black,
  above=5mm of $(lean.north)!0.5!(review.north)$] {SELF-CONTAINED PUBLIC ARTEFACT};
\draw[paper boundary] ($(reason.east)!0.5!(lean.west)+(0,-9mm)$) --
  ($(reason.east)!0.5!(lean.west)+(0,12mm)$);
\end{tikzpicture}
\caption{The end-to-end architecture. The dashed vertical line marks the
release boundary. Everything to its right remains usable without the private
workbench.}
\label{fig:lifecycle}
\end{figure}

The lifecycle has four recurring operations. First, \emph{selection}: choose a
bounded object rather than treating the whole repository as one prompt.
Second, \emph{execution}: run an experiment, edit a proof, or write an
explanation. Third, \emph{validation}: ask the authority appropriate to that
object. Fourth, \emph{binding}: record the result, limitation, and source so a
later agent does not have to rediscover them. The final feedback step turns a
local success or failure into a reusable route, check, or warning when it
genuinely generalises.

The public \path{scripts/proof_cockpit.py} card gives an agent this separation.
It reads checkout and toolchain
identity, corpus scale, registered open propositions, problem obligations, and
workbench sessions from public files.  It imports no private state and adds no
evidence layer to Figure~\ref{fig:architecture}.  Its \texttt{--check} mode
runs the claim-registry, cold-clone, and projection-freshness checks; neither
mode invokes Lean or decides whether prose follows from a theorem.

Once a person has compared a formal theorem with its public wording, a program
can preserve the resulting decision about names, files, wording, and limits.
It cannot decide whether the decision was correct.  The workflow does not
technically force a second independent mathematician: one maintainer can edit
the Lean statement, record, and prose together so that every comparison agrees
with the same mistake.
\subsection{Skill-addressable job lifecycles}
\papersectiontarget{systems-job-lifecycle}

The public clone gives each recurring job one entry skill. A skill owns the
job's starting state, permitted mutations, required evidence, stopping or
re-entry condition, and next owner; it does not acquire the authority of the
objects it coordinates. Figure~\ref{fig:job-lifecycle} records the ordinary
research path.

\begin{figure}[!t]
\centering
\begin{tikzpicture}[
  job/.style={paper stage, text width=2.35cm, minimum height=1.25cm},
  arr/.style={paper arrow}
]
\node[job, fill=MidnightBlue!5] (invoke) at (0,0)
  {{\scriptsize\sffamily\bfseries 1}\\[2pt]\textbf{Invoke}\\[2pt]\scriptsize bounded route};
\node[job, fill=MidnightBlue!5, right=5mm of invoke] (work)
  {{\scriptsize\sffamily\bfseries 2}\\[2pt]\textbf{Work}\\[2pt]\scriptsize continue or suspend};
\node[job, fill=TealBlue!5, right=5mm of work] (validate)
  {{\scriptsize\sffamily\bfseries 3}\\[2pt]\textbf{Validate}\\[2pt]\scriptsize computation or Lean};
\node[job, fill=TealBlue!5, right=5mm of validate] (propagate)
  {{\scriptsize\sffamily\bfseries 4}\\[2pt]\textbf{Propagate}\\[2pt]\scriptsize classify consequences};
\node[job, fill=ForestGreen!6, below=8mm of propagate] (package)
  {{\scriptsize\sffamily\bfseries 5}\\[2pt]\textbf{Package}\\[2pt]\scriptsize evidence and provenance};
\node[job, fill=ForestGreen!6, left=5mm of package] (propose)
  {{\scriptsize\sffamily\bfseries 6}\\[2pt]\textbf{Propose}\\[2pt]\scriptsize pull request or issue};
\node[job, fill=BurntOrange!7, left=5mm of propose] (adopt)
  {{\scriptsize\sffamily\bfseries 7}\\[2pt]\textbf{Review and adopt}\\[2pt]\scriptsize reconcile, accept, credit};
\draw[arr] (invoke) -- (work);
\draw[arr] (work) -- (validate);
\draw[arr] (validate) -- (propagate);
\draw[arr] (propagate) -- (package);
\draw[arr] (package) -- (propose);
\draw[arr] (propose) -- (adopt);
\end{tikzpicture}
\caption{The clone-local job lifecycle. Closure means that the present delta
has evidence and dispositions; it does not mean that the open problem is
solved.}
\label{fig:job-lifecycle}
\end{figure}

The current owners are \texttt{explain-public-system} for reader orientation,
\texttt{run-coupled-research-goals} for discovery--stewardship coordination,
\texttt{mine-open-problem} for invocation and research,
\texttt{lean-concurrent-validation} for bounded Lean checks,
\texttt{propagate-research-consequences} for downstream reconciliation,
\texttt{erdos-research-return} for packaging and maintainer assimilation, and
\texttt{submit-pull-request} for preparing the Git proposal. The separate
\texttt{add-open-problem} lifecycle moves through proposal, incubating formal
lane, and fully indexed public world. \texttt{public-mathematical-writing}
governs a paper projection when propagation reaches exposition. These names
are repository entrypoints rather than a second claim taxonomy.

\section{The private workbench}\label{sec:private}

\subsection{Disk is shared memory}

The private system does not treat a chat transcript as the project. Durable
state lives in files: source code, structured records, append-only events,
generated views, validation receipts, and authored explanations. A conversation
may propose a change, but the file and its governing checker determine whether
the change exists. This makes work resumable across agents and time: the next
agent reads a bounded current-state packet and the relevant sources rather than
trusting a summary of an earlier conversation.

Different file classes carry different authority. Raw operator language
preserves what was asked. Work records say what has been claimed and by whom.
Source files contain implementations and proofs. Receipts record what actually
ran. Authored papers explain; generated indexes help readers navigate. A
generated projection is therefore a map of the territory, not the territory's
source of truth.

At entry, a deterministic router selects a small context packet for the task.
The packet gives candidate objects, authority boundaries, relevant standards,
and legal next actions. This ``small map before large source'' rule reduces
blind search without pretending that a summary is proof.

\subsection{Type A and Type B}

The system uses \emph{Type A} and \emph{Type B} to describe substrate access,
not model quality.

\begin{table}[!ht]
\centering\small
\begin{tabular}{@{}>{\raggedright\arraybackslash}p{0.16\linewidth}
                    >{\raggedright\arraybackslash}p{0.34\linewidth}
                    >{\raggedright\arraybackslash}p{0.39\linewidth}@{}}
\toprule
\papertablehead
Actor & What it can do & Governing limit \\
\midrule
Type A & Inspect live project state; use repository tools; edit claimed files;
run tests; bind receipts and status. & It may change only the substrate and
paths its task authorises, and its claims remain limited by the relevant
validator or human review. \\
\addlinespace
Type B & Reason over a selected packet through an external model, web, API, or
operator-carried exchange; return research, critique, or a candidate. & It has
no direct private-substrate authority. A Type A actor or the operator must
check and apply useful output. \\
\bottomrule
\end{tabular}
\caption{Type A/B is an authority distinction. Either type may be weak or
strong; a delegated tool-using agent is still Type A if it has live substrate
access.}
\label{tab:typeab}
\end{table}

This separation has two benefits. External reasoning can be used without
pretending it observed private state, and local mutation stays attached to
tests, paths, and receipts. It also prevents a common category error: calling
every delegated worker Type B, or treating Type B as a cheaper or less capable
kind of intelligence.

\subsection{Concurrency without shared-state confusion}

Several agents may work at once, but concurrency is admitted only where the
write scopes can be separated. A work item identifies the objective and
expected evidence. Before mutation, an agent claims exact paths for a bounded
lease. Independent claims may proceed concurrently; overlapping claims must be
coordinated or deferred. Fan-out is followed by a fan-in barrier where the
results are compared, validated, and integrated.

Coordination state is written to append-only ledgers and immutable receipts,
not inferred from who last spoke in chat. Generated status pages are read
models over those records. This gives three distinct facts: an agent may be
running, its proposed change may exist, and its change may have passed the
owner's gate. None implies the next. Bounded job counts and focused Lean builds
also prevent concurrency from exhausting the machine or making two builds
corrupt each other's evidence.

\subsection{The control plane is executable}

The router projects a typed option surface over skills, standards, paper
modules, live work, and source artefacts, then
selects a task-conditioned packet before large sources are opened. A cold
agent can move from a one-line flag to a short card, a working context, and
finally the authority-bearing file. Each compression layer carries its source,
safe drill-down, and authority posture. Context selection is therefore an
inspectable operation rather than an invisible choice by the language model.

The work ledger separates past work from present permission. Completed,
reopened, and superseded work are append-only lifecycle events. Live claims are
short leases over exact paths or work objects. Directory claims collide with
their children; expired leases lose authority unconditionally; a crashed
holder leaves an expiry and dirty-handoff record rather than silently
disappearing. A generated cohort view may report activity, collisions, and
unknown scope, but it cannot grant a write. Thus ``an agent mentioned this
file'', ``an agent holds this file'', and ``a change to this file was accepted''
remain different statements.

An additional resident-runtime layer, called metabolism, turns repository
events into bounded maintenance jobs. One daemon owns a SQLite store in
write-ahead-log mode with events, jobs, runs, provider budgets, heartbeats, and
temporal blackboard claims. Hooks, filesystem scans, provider interruptions,
and explicit commands feed the same store. Stable digests collapse duplicate
events; an allowlist limits what may execute; cooldowns smooth provider use;
expired owners are recoverable; and a single-resident guard prevents two
schedulers from fighting over the queue. Status pages are projections over the
store, never repair authority. During one snapshot used for this revision the
status surface correctly reported the daemon as not running while queued jobs
remained visible. The substrate exists and preserves work, but ``always-on
architecture'' is not the same claim as ``currently healthy service''.

The operating loop is deliberately plain: observe the current state, classify
the task, route it to an owner, claim the work, act, validate, record the
result, propagate any reusable lesson, then observe again. ``Continuous'' work
means repeating that loop with explicit checkpoints. It does not mean allowing
an agent to mutate indefinitely without a new receipt.

\subsection{What continuous mathematical work looks like}

A continuous goal is not a request to repeat the same prompt forever. It is a
durable research object with a fixed ultimate question, a mutable frontier,
and an explicit claim ceiling. Each new run resumes from the latest accepted
state: the sources already read, the routes already killed, the computations
already interpreted, the formal obligations still open, and the precise
statement that remains unproved. The agent then selects the highest-value
available transition in that state. It may close a lemma, produce a
counterexample, sharpen a reduction, repair a source claim, or expose a
smaller obstruction. Merely producing more activity is not a transition.

A representative internal trace makes this less abstract. Over fifteen
successive turns, one problem-directed run recorded 313 visible progress
updates and 3,491 command events, with linked workers exploring separate
analytic, computational, and formal lanes. The movement was not a straight
line from prompt to proof. Numerical probes exposed structure; exact arithmetic
replaced promising samples; several attractive strengthenings were killed by
counterexamples; partial inequalities were narrowed to their valid domains;
Lean targets were attempted only after an ordinary mathematical kernel had
stabilised; and successful local results were propagated into the problem
frontier before the next search began. While long exact computations ran, the
controller advanced independent work instead of repeatedly polling them.

That trace also displays why the architecture keeps several evidence layers.
Its compact narrative reported no observed changed paths even though its own
turn summaries referred to landed commits. This is not evidence that nothing
changed, nor that the summaries were correct. It is evidence that a compressed
trace has an observation boundary. Repository state, command results, diffs,
formal builds, and commit objects must settle the question. Likewise, four
turns lacked complete closeout events and only sixteen of forty-two linked
session outcomes appeared in the compact projection. Completeness is therefore
an explicit field, not an impression created by fluent prose.

The example is illustrative rather than a throughput benchmark. It shows the
control shape actually used: search and validation interleave; failed routes
remain informative; concurrency is conditional on separable work; and every
claimed advance must eventually cross from a narrative of activity into an
authority-bearing artefact and receipt. The design aim is continuous
mathematical pressure with discrete, inspectable commitments. A run may stop
because the next step needs an unavailable validator, a human decision, or a
new idea. It must then preserve a useful re-entry point rather than manufacture
busywork or quietly weaken the endpoint. Long-running computation is likewise
not a reason to poll idly when an independent proof, source, exposition, or
adversarial lane remains safe to advance.

\subsection{Coupled continuous goals: discovery and stewardship}
\papersectiontarget{systems-coupled-goals}

One continuous goal is not enough for a moving mathematical corpus. A
problem-directed miner can remain close to the proof frontier, yet that same
proximity makes it a poor sole judge of whether a new lemma is routine,
whether an older theorem is stronger, or where the result belongs in a paper.
The implementation therefore admits two coupled goals with different
responsibilities. The \emph{discovery goal} reads the current problem world,
runs discriminating computations, attempts proofs, records falsifiers, and
returns the smallest stable mathematical delta. The \emph{stewardship goal}
reads that delta against the whole source-current corpus, composes related
declarations into coherent result families, and reconciles their downstream
uses. Neither role receives the other's authority.

\begin{figure}[!t]
\centering
\begin{tikzpicture}[
  role/.style={paper stage, text width=3.05cm, minimum height=1.5cm},
  object/.style={role, fill=TealBlue!6, draw=TealBlue!60!black,
    text width=2.72cm},
  output/.style={role, fill=ForestGreen!6, draw=ForestGreen!55!black,
    text width=3.15cm},
  arr/.style={paper arrow},
  feedback/.style={paper feedback}
]
\node[role, fill=MidnightBlue!5, draw=MidnightBlue!55] (mine) at (0,0)
  {{\scriptsize\sffamily\bfseries DISCOVERY}\\[3pt]\textbf{Test the frontier}\\[2pt]
   \scriptsize compute, prove, or falsify};
\node[object, right=7mm of mine] (delta)
  {{\scriptsize\sffamily\bfseries HANDOFF}\\[3pt]\textbf{Landed delta}\\[2pt]
   \scriptsize theorem, no-go, computation};
\node[role, fill=BurntOrange!7, draw=BurntOrange!65!black,
  right=7mm of delta] (steward)
  {{\scriptsize\sffamily\bfseries STEWARDSHIP}\\[3pt]\textbf{Place the result}\\[2pt]
   \scriptsize appraise, reconcile, explain};
\node[output, below=9mm of delta] (surfaces)
  {{\scriptsize\sffamily\bfseries PUBLIC SURFACE}\\[3pt]
   \textbf{Current mathematical picture}\\[2pt]
   \scriptsize paper order, evidence, frontier};
\draw[arr] (mine) -- (delta);
\draw[arr] (delta) -- (steward);
\draw[arr] (steward) -- (surfaces);
\draw[feedback] (surfaces.west) to[out=180,in=-90]
  node[below left, font=\tiny\sffamily, align=center]
  {ranked next question, missing evidence,\\discriminating test} (mine.south);
\end{tikzpicture}
\caption{The coupled continuous-goal lifecycle. A landed mathematical delta
wakes stewardship; a changed appraisal or consumer gap can change the next
mining target. An unchanged corpus produces no model heartbeat.}
\label{fig:coupled-goals}
\end{figure}

The stewardship pass makes four decisions separately. \emph{Authority}
records what Lean, a computation, a paper argument, or an external source
actually establishes. \emph{Mathematical appraisal} asks about logical reach,
mechanism depth, independence, sharpness, reuse, and the surviving open
boundary. \emph{Exposition placement} determines what leads the abstract,
which theorem receives the longest explanation, and what remains subordinate.
\emph{Work allocation} determines which missing implication or consumer
deserves the next unit of compute or expert attention. A result can be highly
significant but mechanically unready, fully checked but mathematically
routine, or ideal as a worked example without being the corpus's strongest
theorem. No single scalar or status is allowed to decide all four questions.

This is also a second checking layer, but not a second proof authority. The
steward may detect that a purported advance merely restates the open problem,
that several theorem names express one result, that a paper has buried a
stronger theorem, that Comparator exposes a weaker interface, or that a
Palomar packet omits the hard mechanism. It may then narrow, regroup,
reorder, or defer the projection. It may not turn significance, repeated
agreement, or polished exposition into a proof. Lean validation, intended-
meaning review, prior-art assessment, and community acceptance remain
distinct gates.

The clone-local \texttt{run-coupled-research-goals} skill operationalises this
division by invoking \texttt{mine-open-problem} and
\texttt{propagate-research-consequences} as separate jobs with source-pinned
hand-offs. A single agent may execute the two jobs sequentially; two tasks,
machines, models, or people may also hold the roles. The separation is logical,
not a requirement to buy two computers.

The coupling is event-driven. A newly landed theorem, counterexample,
authority change, paper correction, or external review outcome can wake the
stewardship goal. Its output is a source-pinned appraisal and a changed
consumer or frontier disposition, not a recurring status report. That output
may in turn wake the discovery goal with a stronger target, a missing
hypothesis, a counterexample request, or the cheapest discriminating
computation. If neither the mathematical frontier nor a downstream consumer
has changed, both goals yield rather than polling one another.

Before either role begins, \texttt{explain-public-system} provides a separate
cold-clone orientation job. An agent can read the public corpus and companion
papers, adapt the explanation to a lay reader, mathematician, formaliser,
compute contributor, reviewer, or infrastructure contributor, and point to the
exact evidence behind its account. The human therefore need not learn the file
layout before asking a useful question. Explanation remains a projection: it
cannot promote its own summary into proof or acceptance.

\section{The mathematical reasoning loop}\label{sec:mathloop}
\papersectiontarget{systems-mathloop}

The mathematical lane begins before Lean. A problem is first stated in ordinary
mathematical language, with its known status and source. Candidate mechanisms
are then generated and ranked. Rank is based on mathematical signal---endpoint
proximity, depth of mechanism, independence from other routes, usefulness for
future work, and overclaim risk---rather than theorem count or ease of
formalisation. Attention is deliberately unequal: a decisive obstruction or a
close conditional reduction may deserve more work than many routine lemmas.

\subsection{Experiments are route selectors}

Computation serves three useful roles. It can find a counterexample, measure a
finite pattern, or test whether a proposed mechanism is plausible enough to
formalise. Each experiment records its inputs, code, finite domain, output, and
interpretation. The interpretation must state what the experiment cannot show.

A finite computation becomes formal evidence only through an explicit bridge.
For example, Lean may evaluate a finite proposition with \texttt{decide} and
check the resulting proof term. Even then, the conclusion remains finite. A
pattern observed for many inputs does not become an ``all inputs'' theorem, and
a successful numerical approximation does not become an exact equality. Failed
experiments are kept when they prune a natural route; otherwise later agents
pay to repeat the same mistake.

\paragraph{Failure modes are mathematical outputs.}
The system distinguishes three kinds of negative evidence. A failed agent
attempt says only that one search path did not close. A counterexample can
refute a conjecture on its exact domain. A Lean no-go theorem can rule out a
whole strategy class under explicit hypotheses. The last two may be as
valuable as a positive lemma: they prevent repeated work and reveal which new
ingredient an endpoint requires. This corpus therefore keeps deep
problem-specific reasoning surfaces, proof gaps, coefficient-only and
fixed-precision obstructions, and corrected statements beside successful
theorems. Every no-go keeps its scope visible; failure of one mechanism is not
failure of every possible proof.

\subsection{Three oracles, not one}

Experiments are route selectors; they do not become a fourth oracle.
The workflow separates three answers that are often compressed into the word
``verified''.

\begin{enumerate}
\item A \emph{status oracle} says what a public registry or the literature
      currently reports about the problem.
\item A \emph{formal oracle} says whether the pinned Lean environment accepts
      this exact statement and proof under the permitted axioms.
\item A \emph{research-validity judgement} asks whether the formal statement
      matches the intended problem and whether the result matters in context.
\end{enumerate}

The second can be mechanical. The first can become stale. The third remains a
mathematical and scholarly judgement. A successful Lean build without statement
reconciliation is therefore not the end of the reasoning process.

The learning loop grows the corpus it reasons over. An agent first navigates
the existing Lean and semantic graph, then uses proof search and bounded
computation to probe promising gaps and patterns that may be too distributed
for an unaided manual scan. The checking pass records the kernel result,
reconciles formal and informal statements, and packages either a theorem, a
counterexample, or a diagnosed failure. Accepted nodes and edges improve the
next search; reusable process lessons improve heuristics and context, never
mathematical truth. When a mechanism comes from the literature, its annex
record, locator, and public citation preserve the original authorship and
claim scope while Lean checks only the new formal statement. Thus the cycle
is: ground sources, navigate, compute, conjecture, formalise, interpret,
publish, and return the enlarged graph to the next iteration.

\subsection{From local progress to reusable mathematics}

Problem-sized worlds are deliberately local, but their useful by-products
need not remain local. A lemma discovered while attacking one Erd\H{o}s problem
may encode a comparison principle, finite combinatorial device, analytic
estimate, or Lean interface that belongs in a wider mathematical library. The
system therefore distinguishes \emph{process up-propagation}---a lesson that
improves how later agents work---from \emph{mathematical
canonicalisation}---the conversion of a local result into a reusable theorem.

Canonicalisation is a second research act, not a change of filename. It asks
which hypotheses the proof really uses, removes problem-specific coordinates,
searches for prior art and existing library interfaces, and states a natural
general theorem for which the motivating local result is an explicit
specialisation. The candidate is then attacked again: dropped hypotheses need
counterexample tests, the general proof needs its own exact check, and at least
one genuine additional consumer or explanatory use should be identified.
Broader syntax without broader mathematical use is not yet a canonical
result.

This yields three statuses that must not be collapsed: checked inside one
problem world; proposed as reusable mathematics; and suitable for an external
shared library. The present system can preserve the first, help construct and
test the second, and prepare evidence for the third. It cannot grant the third
status to itself. Mathlib is open to new contributors, but its current guide
sets high standards for generality, integration, maintainability, style,
documentation, and responsible disclosed use of AI \cite{mathlib}. A subject
and Lean expert must decide whether a candidate belongs there, reshape and
review it as needed, and take responsibility for any upstream discussion or
pull request. The originating route should remain visible in provenance; the
expert should receive explicit credit for the generalisation, library design,
formalisation, review, and stewardship actually supplied. The
\papersectionlink{open-source-mathematics-strategy.pdf}%
{strategy-local-to-general}{open-source strategy} gives the corresponding
participation and credit protocol.

\subsection{Problem-sized Lean worlds and bounded theorem neighbourhoods}

Large formal corpora need several maps because ``everything is indexed'' is
not the same as ``everything is understood.'' The declaration atlas answers
``where is it?''; exact imports and dependencies answer ``what does it formally
use?''; authored semantic relations answer ``what role does it play?''; and the
claim graph answers ``what may the project publicly say?'' Generated or inferred
edges can select evidence, but only source-specific gates can supply proof,
interpretation, or publication authority.

The scale is already problem-sized. At the August 2026 snapshot the attached
public corpus described 1,024 Lean modules and 153,396 declarations. The much
deeper private record for Problem 1041 contained 177 exact results, seven open
producers, 72 negative results, 134 Lean modules, and 335 executable experiment
programs. Its public research-corpus manifest alone recorded 685 files and 35
selected strongest results. These are dated inventory facts, not a throughput
benchmark; generated certificate families also make raw theorem counts a poor
measure of mathematical value.

Depth is not readiness. The canonical public 1041 library currently has two
integrated Lean modules, while its separately governed public research export
has 125 Lean sources and its private world is larger again. The public route
memory still records no reviewed 1041 claim family. Across the whole semantic
corpus, structural linkage is broad, but only 25 statement nodes and eight
relations carry digest-bound semantic-review receipts at this snapshot. These
differences are intentional status boundaries: a discoverable declaration is
not thereby understood, integrated, reviewed, or publishable.

An agent does not receive this world as one prompt. First the problem cockpit
fixes the endpoint, status, claim ceiling, and canonical frontier. It keeps open
producers distinct from the target so a promising route cannot silently replace
the problem. Next a bounded neighbourhood selects the strongest relevant
results, exact source coordinates, imports, consumers, alternative formulations,
experiments, counterexamples, no-gos, and literature. It labels every item by
evidence class and emits an omission receipt with an expansion route. One 1041
packet exposed only four of 134 Lean modules, four of 335 experiments, and four
of 55 no-go notes; designed omission, not exhaustive loading, made the context
usable. When an exact path is claimed, a local connection card can put its
prerequisites, sibling mechanisms, consumers, falsifiers, and validation target
before broader retrieval.

\subsection{Comprehension before the mathematics}
\papersectiontarget{systems-comprehension}

The neighbourhood is compiled, not remembered. Before attempting mathematics an
agent asks a working-memory compiler for a packet conditioned on its query, and
that packet, rather than the repository, is what it reasons over.

Federation works by identity. Each corpus keeps its own authority and nothing
is copied into a central index. The compiler holds a compact descriptor and a
content fingerprint for the private projection and for every attached public
checkout, and it opens an attached exhaustive atlas only after those identities
reconcile and the bounded route has failed to answer. It never rebuilds the
expensive Lean microcosm. Every packet names the corpora it consulted and the
digests at which it saw them.

The packet leads with the endpoint. A problem cockpit fixes the target
statement, its status, the claim ceiling, and the current claim frontier,
labelling each frontier claim with an evidence class and a logical altitude
that separates an equivalence from a reduction, a necessary condition, a finite
exclusion, and a no-go. Beneath it a mechanism landscape keeps proved or
kernel-checked results apart from open producers and from routes a
counterexample has closed. Retrieval moves between global frontier structure,
concept neighbourhoods, and premise or obligation ancestry, and exact imports,
authored arguments, generated navigation, lexical references, and semantic
inference remain visibly different edge classes.

Above that material sits an inference workspace: an exact commitment plane for
source coordinates and statements, a strategy plane whose unit of branching is
a coherent premise group, and an adversary plane for the falsifiers a branch
must survive. Its commitments are invalidated whenever a corpus fingerprint
moves. It states its own boundary in every packet, that only Lean, an exact
counterexample, or an owner verifier receipt may change the status of a claim.
The workspace plans inference and concludes nothing.

Two refusals are as important as the retrieval. A query naming no theorem,
declaration, or claim is reported as unanchored, and its consumer is told to
name one before treating any route as target-specific, so a vague question does
not receive a specific frontier. A packet over its requested context budget
says so and exits non-zero instead of truncating quietly: it reports the
requested and estimated token counts, names the planes it dropped, and gives
the command that restores them. An omitted adversary plane appears as an
omission with an expansion route, never as an absence of falsifiers.

This is orientation, not understanding. A packet can be well formed over
corpora it identified correctly and still select a weak mechanism, omit the
governing obstruction, or carry an authored interpretation that is
mathematically wrong. The supported claim is narrower: an agent entering a
corpus of this depth receives a bounded starting state that names its sources
by digest, labels its evidence, and declares its own omissions, and no part of
that state carries proof authority. Whether agents working from such a packet
produce better mathematics than agents without one is not measured here.

\subsection{Consequence propagation}

When a declaration lands, the direction reverses. A consequence mapper begins
at the exact changed object and enumerates reverse imports, claim records,
validators, experiments, papers, and open obligations that may now be stale.
A semantic second pass must choose: update now, verify unchanged, defer with a
reason, or mark outside scope. Empty lexical search is not evidence of no
consequence, and no projection may bulk-strengthen a family of claims. A formal
evidence cell carries the result, its explicit non-claim, evidence class,
source and receipt links, and next unresolved obligation through this fan-out.
The public \texttt{propagate-research-consequences} skill implements the same
discipline without depending on the private workbench. For work returned from
an older clone, it runs once against the contributor's recorded starting
commit and again after the reviewed change has been reconciled with current
main. The original delta and any conflict-resolution delta remain separate
evidence and receive separate attribution.

This same architecture distinguishes mathematical propagation from process
propagation. A theorem, counterexample, or no-go changes the mathematical graph
only through its own verifier. A reusable failure in navigation, scheduling,
experimentation, or validation enters a guarded packet containing the local
case, failure class, evidence, sibling scan, proposed owner, overgeneralisation
guard, validation, and stop condition. If accepted, it may change a route,
skill, check, or standard for later agents. The mechanism is implemented and
has worked examples, but there is not yet a complete measure of what fraction
of useful failures propagate. The supported claim is that failures \emph{can}
alter the durable control plane without altering mathematical truth.

\section{The public Lean repository}\label{sec:public}
\papersectiontarget{systems-public}

The public checkout begins at a deliberate boundary. It contains every file
needed to inspect its claims and replay its checks. It does not call back into
the private workbench, and no public theorem depends on an unpublished private
lemma. The larger workflow is provenance: it explains production, not truth.

The repository has two Lean roots. The reviewed root contains the established
249/257 publication lane. The problem-owned expansion root contains work on six
additional Erd\H{o}s problems and unpromoted lanes for 249 and 257. Both roots
are checked by the same Lean kernel, but kernel acceptance does not automatically
promote a declaration into a reviewed public claim. Promotion is a separate
change to the claim record and exposition.

The main public sources have deliberately separate jobs:

\begin{table}[H]
\centering\small
\begin{tabular}{@{}>{\raggedright\arraybackslash}p{0.23\linewidth}
                    >{\raggedright\arraybackslash}p{0.31\linewidth}
                    >{\raggedright\arraybackslash}p{0.34\linewidth}@{}}
\toprule
\papertablehead
Surface & Authority & It does not establish \\
\midrule
Lean source and pinned toolchain & Exact formal statements and proofs accepted
by the kernel. & Intended meaning, novelty, significance, or faithful prose. \\
\addlinespace
Reviewed claim record & Approved wording, status, evidence, bounded domain,
and adjacent open statement for selected claims. & Correctness or completeness
of the human review. \\
\addlinespace
Papers and guides & Human explanation and reading order within the recorded
claim ceiling. & New formal authority. \\
\addlinespace
Generated maps and query tools & Bounded navigation across declarations,
problems, graphs, papers, and claims. & Proof or permission to strengthen a
claim. \\
\addlinespace
Release checks and continuous integration & That configured identities,
relationships, generated files, licences, and negative tests pass on a named
revision. & Understanding of unrestricted prose or independent mathematical
approval. \\
\bottomrule
\end{tabular}
\caption{The public authority split.}
\label{tab:authority}
\end{table}

The companion public Plectis repository has a different role. It publishes
runnable, bounded mechanism slices and receipts from the wider system. The Lean
repository publishes a mathematical corpus. Neither public repository inherits
authority from the other, and neither is a public mirror of the private root.

\section{Comparator, Palomar, and publication}\label{sec:assurance}

\subsection{Comparator: an exact-statement firewall}

Comparator protects a narrow but important boundary. Selected propositions are
declared again in a challenge module without their proofs. A solution module
must provide terms of those exact types. The configuration pins the permitted
axioms, the modules, and a runtime receipt. A named altered statement must fail,
which checks that the harness has not become vacuous.

Comparator therefore answers: ``Does the proof-bearing corpus still implement
this separately stated Lean interface under this axiom budget?'' It does not
answer whether the interface is the right translation of an informal problem,
whether the theorem is new or important, or whether an independent human has
reviewed the proof. ``Comparator-checked'' is accurate; ``independently
verified'' is not.

\subsection{Review selection, and what the Palomar registry is not}

Comparator's roster is an evidence inventory, not a ranking. A local
review-selection layer adds the missing editorial step. It groups declarations
into result families, ranks them by mathematical signal, keeps the hard
mechanism and surviving boundary adjacent, and selects a small review
portfolio. A qualification check can say that the local packet satisfies its
structural requirements. Its records sit under a Palomar-named showcase file
because selected families may then be submitted to the external Palomar
registry of Lean-verified mathematics. Palomar's own documentation describes
mechanical proof checks and automated editorial filtering; it performs no human
peer review and does not rank significance. Local qualification, registry
inclusion, and human mathematical review are different events. None, by
itself, establishes broad acceptance or significance within the mathematical
community.

This separation matters under proof abundance. Counting theorems rewards
generated volume and routine closure. The local selection instead asks which result changes
the mathematical picture, which conditional route is closest to an endpoint,
which obstruction prevents wasted work, and which explanation will help an
expert assess the claim. The ranking may guide attention; it cannot alter proof
status.

\subsection{Publication and propagation}

After formal proof and statement reconciliation, a result may enter an authored
paper, a claim record, a Comparator packet, and a local review unit. Each is a
separate projection with a separate ceiling. A public release is made only
after the Lean build and the release-surface checks pass, and publication
remains a human action.

The return path is equally important. A proof failure may improve a
formalisation heuristic. A counterexample may close a family of tempting
routes. A reviewer objection may require a narrower public claim. A release
drift may create a stronger check. Those lessons propagate to the smallest
durable owner---a theorem, experiment record, skill, standard, route, or claim
boundary. They do not rewrite raw intent, generated views, or mathematical
status by implication.

\subsection{From formal checking to mathematical use}

Tao argues that AI mathematics should be judged along a chain: proof
generation, verification, exposition, publication and community digestion,
then eventual canonicalisation \cite{taoai}. This architecture is a concrete
attempt to build for that whole chain. The private reasoning and experiment
loops support generation. Lean checks formal proofs, while Comparator protects
selected exact interfaces. The problem papers, reasoning surfaces, graphs, and
claim records support exposition. Local review selection triages scarce expert attention, and
the contribution and release paths support attributable publication. Digestion,
acceptance, and canonicalisation remain achievements of the mathematical
community, never local status fields.

This mapping also explains the unusual emphasis on failure. Tao warns that
AI-polished exposition can erase the natural friction that tells a reader
where the real mathematical difficulty lies. The reasoning surfaces and no-go
graph preserve that friction deliberately: false starts, corrected claims,
missing producers, and formal obstructions remain adjacent to the successful
theorems. The purpose is not to display internal noise. It is to transmit the
shape of the problem so another mathematician can understand why the surviving
boundary is hard and where a genuinely new idea must enter.

Paper authoring itself participates in this loop. As the mathematical and
control graphs change, agents assemble and revise problem papers, architecture
papers, reviewer cards, and claim records from the source-current evidence.
Writing is therefore an active digestion and interpretability pass in Tao's
sense: a missing explanation, unstable term, hidden quantifier change, or
unregistered limitation discovered while writing becomes a typed defect that
can return to the theorem, semantic relation, claim record, route, or validator.
This manuscript is a reflexive example: it was rewritten by the system while
agents inspected the mechanisms it describes, then compiled and checked as a
public artefact. That reflexivity is provenance, not validation. A system does
not independently review itself merely by producing a clear account of itself.

The remaining design choices follow the same programme. Tool use is disclosed;
models are not listed as authors; human contributors retain differentiated
credit; references and source paths are inspectable; automatic checks filter
work without impersonating peer review; and theorem counts do not determine
value. The repository turns those recommendations into an implementation
experiment: it asks how a small project can
prepare for proof abundance without confusing abundant output with collective
mathematical progress.

\section{One complete boundary: finite is not unbounded}\label{sec:example}

Erd\H{o}s Problem 249 asks whether
\[
S=\sum_{n\ge1}\frac{\varphi(n)}{2^n}
\]
is irrational, where \(\varphi(n)\) is Euler's totient function
\cite{erdosgraham}. The development reduces irrationality to a family of exact
non-integrality certificates. Along the diagonal
\(H_t=\operatorname{lcm}(1,\ldots,t)\), write \(\mathrm{Cert}(t)\) for the
existence of the required finite certificate.

Lean checks the finite theorem
\[
\forall t\le82,\quad \mathrm{Cert}(t).
\]
The endpoint needs an unbounded supply:
\[
\forall T,\quad \exists t>T,\quad \mathrm{Cert}(t).
\]
The development proves that the unbounded statement is equivalent to the
irrationality claim. It does not prove the missing implication from the finite
range to the unbounded statement. No matter how large a fixed checked bound is,
a larger cutoff exists.

This example passes through every layer. Computation constructs finite data.
Lean checks the certificate predicate and the finite theorem. The formal graph
records its dependencies. The semantic graph identifies its role. The claim
record states the finite range and names the unbounded requirement as open.
The paper explains why the quantifiers differ. Comparator checks selected exact
interfaces. The local review selection may rank the result family for review. The release checker
requires the limitation to remain visible.

The last check was added because the boundary once escaped. A historical README
edit changed a clause saying that the finite cases did \emph{not} supply the
open requirement into one saying that they completed it. Lean was untouched,
and the release checker passed because that prose relationship had not been
registered. After the relationship was added to the claim record, a deliberately
false copy was rejected. This is evidence for one failure and repair. It is not
a detection rate, a proof that every claim-bearing sentence is registered, or
an evaluation of mathematical review quality.

The historical study was narrower than a benchmark. In that study, nine of the
ten edits were rejected. One escaped because the relationship had not been registered. The
original run logs were not retained. The edits were authored by the checker's
author. After the relationship was registered, only the escaped edit was
reconstructed against the repaired checklist. The other nine edits were not
rerun against the extended checklist. The finite theorem makes no
\(t=83\) or cofinal claim. These limitations are part of the evidence, not
footnotes to be discarded after the repair.
The evidence marks a coverage boundary, not a reliability score.
The post-repair witness accepts the current README and
rejects a test copy containing the false clause.

The same architecture handles other shapes of progress. Problem 257 separates
a theorem for full-support representations from a stronger arbitrary-support
statement that remains open. Problem 68 has an exact carry-based equivalence
but no theorem producing the required carries. Problem 251 has an exact series
reformulation rather than an irrationality proof, and Problem 243 records a
conditional recovery theorem whose premises remain part of the public claim.
A counterexample or corrected statement, as in the 1041 lane, is also a
first-class result. The publication rule is identical in every case: preserve
the quantifiers, premises, and nearest stronger open statement.

\section{Inspection routes}\label{sec:routes}

A reader can inspect the public system through short, question-shaped routes.
The labels below link to repository paths; the paper does not print the full
URLs.

\begin{table}[H]
\centering\small
\begin{tabular}{@{}>{\raggedright\arraybackslash}p{0.34\linewidth}
                    >{\raggedright\arraybackslash}p{0.54\linewidth}@{}}
\toprule
\papertablehead
Question & Start here \\
\midrule
How does the public repository fit together? &
\repolink{ARCHITECTURE.md}{ARCHITECTURE.md} \\
What may the project say, and what remains open? &
\repolink{docs/claims.json}{docs/claims.json} and
\repolink{docs/methodology.json}{docs/methodology.json} \\
Where is the checked mathematics? &
\repolink{Erdos249257.lean}{Erdos249257.lean} and
\repolink{ErdosProblems.lean}{ErdosProblems.lean} \\
How can I navigate without reading the whole corpus? &
\repolink{docs/ORIENTATION.md}{docs/ORIENTATION.md} and
\repolink{scripts/query_corpus.py}{scripts/query\_corpus.py} \\
What exactly does Comparator check? &
\repolink{docs/EXTERNAL_VERIFICATION.md}{docs/EXTERNAL\_VERIFICATION.md} and
\repolink{verification/comparator.json}{verification/comparator.json} \\
What does the local review selection qualify? &
\repolink{docs/PALOMAR_QUALIFICATION.md}{docs/PALOMAR\_QUALIFICATION.md} and
\repolink{docs/PALOMAR_RESULT_SHOWCASE.json}{docs/PALOMAR\_RESULT\_SHOWCASE.json} \\
Which papers exist and what question does each answer? &
\repolink{docs/papers/README.md}{docs/papers/README.md} \\
Which checks gate a release? &
\repolink{scripts/check_release.py}{scripts/check\_release.py} and
\repolink{.github/workflows/lean.yml}{.github/workflows/lean.yml} \\
How can work return with public credit? &
\repolink{CONTRIBUTING.md}{CONTRIBUTING.md} and
\repolink{docs/research-commons/CONTRIBUTIONS.md}{research-commons/CONTRIBUTIONS.md} \\
\bottomrule
\end{tabular}
\caption{Compact public inspection routes. These are entry points, not a new
authority layer.}
\label{tab:routes}
\end{table}

\section{What can be trusted}\label{sec:trust}
\papersectiontarget{systems-trust}

The architecture supports a chain of narrow conclusions:

\begin{enumerate}
\item a recorded experiment ran on its stated finite inputs;
\item a pinned Lean kernel accepted its stated formal source;
\item a maintainer approved a selected interpretation and public boundary;
\item Comparator matched selected proof-bearing declarations to separately
      stated interfaces and an axiom budget;
\item the local review-selection checks validated the structure of a selected
      review packet;
\item the release program found every configured public relationship intact on
      the named revision.
\end{enumerate}

Joining these receipts improves traceability, but it does not create a seventh,
stronger oracle. In particular, the architecture does not establish that the
formalisation is the uniquely right one, the review is independent, the result
is novel or significant, the registered boundary is complete, or an open
Erd\H{o}s problem is solved. A coordinated but mistaken edit to source, record,
and prose can still make every structural comparison agree. Unregistered prose
can still overclaim. Literature status can still change.

The system reduces these risks by keeping sources plural, boundaries explicit,
and negative tests close to high-risk claims. It does not claim to eliminate
human judgement. It does not technically force a second independent
mathematician to review the result. This is the same limit encountered in
assurance cases: a formal notation can record an asserted support relationship
without proving that the evidence is true or sufficient for the top claim
\cite{gsn}.

A reader can keep those limits straight by asking, for any capability named
in this paper, which of four kinds of support stands behind it.

\begin{center}
\small
\begin{tabular}{p{3.1cm}p{5.2cm}p{5.2cm}}
\toprule
Kind of support & Example & What may be said \\
\midrule
Executable check & A session or return validated against required fields and source identity; a Comparator interface match & The program checks the stated conditions on the stated inputs \\
Agent workflow instruction & The discovery and stewardship responsibilities in a public skill & The workflow directs the agent; it does not show that every agent complied \\
Mathematical judgement & Whether a result expresses the intended mathematics or improves on prior work & A named assessment is required; agreement between source, record, and prose is not one \\
External outcome & Another researcher reproduces, adopts, corrects, or builds on the work & A recorded external event is required; the architecture implies none \\
\bottomrule
\end{tabular}
\end{center}

\section{Scaling from one clone to a search network}\label{sec:scaling}
\papersectiontarget{systems-scaling}

\subsection{Clone, attach compute, and mine bounded results}

The public repository is more than a static archive. A fresh clone fixes a
Lean toolchain, exposes the supported roots, names reviewed claims and open
boundaries, and provides query and release programs. An external researcher
can attach any agent runner to those public interfaces, launch independent
workers on disjoint theorem or experiment lanes, and submit candidate changes
through the same proof and publication gates. More compute increases the
number and diversity of attempts; it does not relax the meaning of a passing
receipt.

This suggests a ``result mining'' mode. Workers select bounded frontier
objects, explore proofs, counterexamples, computations, or literature, and
return exact artefacts. Fan-in then checks Lean, reconciles statements, records
negative evidence, and ranks the mathematically strongest surviving result
families. The useful unit is not a model answer or a theorem count. It is a
result packet with provenance, scope, formal status, mathematical role, and a
nearest stronger open statement. Local review selection can allocate scarce
expert attention to those packets; Comparator can protect selected exact
interfaces; neither is replaced by scale.

At scale, fan-in is a persistent stewardship role rather than an end-of-run
cleanup task. It consumes each stable delta, compares it with the complete
candidate universe, updates the paper hierarchy and assurance interfaces when
warranted, and returns a newly ranked frontier to the mining fleet. Several
miners may therefore share one stewardship goal, and one steward may consume
results from several problems, provided every write and mathematical object
still has one owner. The steward is allowed to say that a large run produced
no new mathematical family, or that a short negative result should outrank
many formal lemmas. This prevents compute volume and theorem count from
becoming accidental editorial policy.

Fleet size and validation capacity are deliberately different variables. In
principle, many logical agents can search disjoint mathematical neighbourhoods;
in practice, the controller admits work according to path conflicts, memory,
CPU, disk pressure, and the expected value of the lane. Cheap reading and
computation may fan out widely while mutation remains path-leased and heavy
Lean work is narrower. This prevents ``more agents'' from meaning ``more
processes contending for the same compiler and checkout.''

Lean validation uses a semantic single-flight queue. A request is identified
by the reachable source, logical targets, toolchain, dependency lock, and
authority mode. Its command string and working directory do not determine
request identity.
Equivalent requests share one owner and may reuse a completed receipt; changed
inputs require a fresh build. Distinct requests that would still compete for
the same host-wide Mathlib resource are serialized by a shared conflict key.
An interactive agent that cannot acquire the slot receives a typed deferral,
not a false theorem failure: it can advance result-independent work while a
detached process owns the queued build. Thus the fleet may grow much wider
than the expensive validator without corrupting the evidence or wasting every
agent as a queue supervisor.

There are consequently four separate scaling limits: logical search width,
safe concurrent mutation, mechanical validation throughput, and scarce expert
review. Path leases govern the second, the Lean queue governs the third, and
review-selection ranking governs the fourth. Fan-in joins their outputs only after
the appropriate receipts exist. This separation is what allows additional
models and compute to enlarge exploration without silently enlarging any
model's authority.

The return path is already public. A person can fork or clone the repository,
work from a named public commit, and open a pull request or research-progress
issue. A structured return can preserve the contributor, collaborators, tool
operator, disclosed model systems, starting commit, evidence, affected result,
and surviving limitation as separate fields. Only an accepted receipt enters
the generated contribution views. Acceptance, mathematical claim status, and
release inclusion remain separate decisions, so a credited counterexample or
failed route need not be mislabelled as a theorem. Later corrections append a
new history rather than erasing the earlier contributor. Attribution and pull
requests are standard practice; the architectural role here is to carry credit
and provenance through fan-in without allowing either to strengthen the claim.

The named starting commit remains useful even when main advances. Git retains
the common ancestor needed to inspect the contributor's original delta. A
maintainer can replay that state, reconcile the reviewed substance with the
current tree, and rerun current validation and consequence propagation. A
material conflict resolution is a new integration contribution; it does not
rewrite the origin of the earlier work.

The exact contributor-facing sequence is specified in the
\papersectionlink{open-source-mathematics-strategy.pdf}%
{strategy-protocol}{companion contribution protocol}; its credit model is
separated in the
\papersectionlink{open-source-mathematics-strategy.pdf}%
{strategy-credit}{provenance and credit section}. The navigation assumptions
made here are tested by the
\papersectionlink{cold-clone-to-proof-receipt.pdf}%
{cold-clone-problem}{cold-clone case study}.

The current public clone supplies the mathematical corpus and its gates. The
full private orchestration layer is not yet distributed as a turnkey public
service, so mass multi-provider mining is a design target rather than a
reported benchmark. The important point is that the public interfaces do not
depend on one private scheduler. A laboratory with many frontier models can
replace the producer while retaining the same evidence boundary.

No outside contributor had completed this path by 31 August 2026. The
cross-paper links, clone-local skills, and pull-request route are therefore
implemented prototype interfaces whose usability remains an external test.

\subsection{The no-go graph as a new mathematical object}

A mature corpus should not represent progress as a list of proved theorems.
Its graph should contain at least the following node classes: endpoint
problems; conjectures and reformulations; proposed mechanisms; computations;
formal lemmas; counterexamples; no-go theorems; reviewed claims; and unresolved
obligations. Edges should distinguish implication, equivalence, dependency,
refutation, obstruction of a strategy, weakening, generalisation, experimental
support, formalisation, and publication. A no-go then occupies a precise place:
it closes one region of mechanism space while exposing the boundary around
nearby variants that remain viable.

As this negative and positive graph grows, several new uses become testable.
A model may navigate around already closed strategy families, identify a small
cut of missing premises between the current corpus and an endpoint, or search
for analogies between obstruction patterns in different problems. Training
data can be formed from proof/counterexample/no-go triples or from contrastive
pairs consisting of a tempting argument and the exact theorem explaining its
failure. Graph-conditioned models might propose the next useful lemma from a
frontier neighbourhood rather than from the theorem statement alone. These
are hypotheses about future systems, not results established by this paper.

Negative structure also creates risks. An over-broad obstruction edge can hide
a viable variant; repeated failed attempts can teach stylistic avoidance rather
than mathematics; and model-generated semantic edges can disagree with the
formal graph. Training snapshots therefore need immutable generations,
source-level provenance, explicit edge authority, and held-out evaluation.
Lean can verify formal nodes and some edges, but expert judgement remains
necessary for semantic roles and for deciding whether the graph captures the
important mathematical space.

\subsection{An experimental agenda}

The next evaluation should ask more than how many theorems additional compute
produces. It should compare graph-aware and graph-blind agents on frontier
selection, repeated-dead-end rate, time to a useful obstruction, premise
discovery, proof completion, and calibration of public claims. It should test
transfer between unrelated mathematical domains as well as between neighbouring
Erd\H{o}s problems. Independent teams should replay result packets and review
whether no-go edges are scoped correctly. The central scaling question is
whether an increasingly rich map of what works and what cannot work makes
future reasoning more efficient and more original, or merely makes the system
more confident about the territory it already knows.

\section{Relation to other approaches}\label{sec:related}
\small

\paragraph{Theorem-proving agents.}
This architecture is not a new proof-search algorithm. LeanDojo and Pantograph
provide programmatic Lean environments and proof-state interaction
\cite{leandojo,pantograph}. OpenProver already combines a planner, parallel
workers, independent verifiers, compact working memory, a larger repository,
and Lean feedback \cite{openprover}. Agent Hunt already studies concurrent
formalisation with locks, bounties, guarded ownership, and collaborative
agents \cite{agenthunt}. DreamProver learns a compact reusable lemma library
through wake--sleep cycles \cite{dreamprover}. These are baselines for proof
interaction, parallel search, and learned proof memory. Here those mechanisms
sit upstream of a different question: after a proof is found, what exactly may
move into a reviewed public claim?

\paragraph{Graphs and checked exposition.}
Proof blueprints connect informal proof plans to named Lean declarations.
\texttt{leanblueprint} checks that author-supplied declaration names exist,
while LeanArchitect infers formal dependencies and unfinished-proof status and
exports synchronised blueprint material \cite{leanblueprint,leanarchitect}.
The graph layers here have a related navigational role, but the public-claim
record begins after a result has been selected and asks which wording and open
boundary were reviewed.

Semantic-audit systems address a neighbouring problem. Lean Atlas narrows the
declarations a person must inspect for chosen theorem statements, conditional
on the semantic correctness of the returned set and trusted base
\cite{leanatlas}. EconCSLib uses models to translate and compare formal and
informal statements, with saved human judgements for a subset \cite{econcs}.
The present architecture uses models throughout production but assigns none of
them final semantic authority. Their output becomes evidence or a candidate
until a source-specific gate accepts it.

Checked-document and traceability systems make heterogeneous relationships
explicit. Isabelle/DOF places formal and informal material in a typed checked
document \cite{isadof}; requirements traceability follows commitments across
development and revision \cite{gotel}. This repository keeps prose unrestricted
and records selected public boundaries separately. That choice makes adoption
simple, but leaves unregistered prose outside the checker.

\paragraph{Auditable scientific agents.}
HEP makes hypotheses, evidence, belief updates, lineage, and resolution states
explicit in an append-only registry \cite{hep}. Symposium proposes immutable
community publication histories with attributable artefacts, declared
evidence, assumptions, and purpose-sensitive arguments \cite{symposium}.
EurekAgent treats permissions, artefacts, budgets, and human supervision as
first-class parts of an agent environment \cite{eurekagent}. Persistent records
and environment design are therefore not unique to this system. The narrower
distinction here is between authorities that must not be collapsed:
experimental evidence, kernel acceptance, reviewed interpretation,
exact-statement assurance, editorial selection, and public release.

Mutation testing asks whether a test distinguishes a seeded fault from the
original program \cite{demillo,jiaharman}. The deliberately false boundary in
the worked example follows that idea. Because the examples were selected by
the system's author and were not run as a controlled external evaluation, they
show that particular checks can fail; they do not measure overall adequacy.

\paragraph{Contribution and limits.}
The contribution is architectural composition, not a claim that each component
was invented here and not evidence of better proof-search performance. The
system joins a private authority-aware workbench to a self-contained public
Lean repository, then keeps formal dependencies, semantic interpretation,
reviewed claims, Comparator interfaces, review selection, and release
relationships visibly separate. It also carries deeply investigated failure
modes and formal no-gos through the same public path as positive theorems,
without discarding the stronger endpoint that survives. Failure logging and
hypothesis refutation have precedents; among the systems compared here, the
distinctive contribution is their integration with formal obstruction results,
claim ceilings, and publication assurance. The evidence is a detailed design
and worked reconstruction from one evolving corpus. It is not a controlled
comparison, an independent audit, or a general reliability estimate.

\paragraph{Scaling beyond this corpus.}
The design is not tied to Erd\H{o}s problems. A new domain needs stable result
identities, source and status records, a formal checker where one exists, and
an explicit public-claim boundary; it need not use the same mathematics or
even Lean. Larger deployments should federate independently owned corpora
rather than build one authority database, separate producer and reviewer
teams, benchmark each boundary independently, and preserve immutable result
generations. Proof-search layers can adopt distributed task markets,
persistent lemma learning, and richer human steering, while publication layers
can add external replay, signed review receipts, and cross-project claim
links. None of these extensions should allow learned process memory or a graph
edge to change mathematical truth by itself.
\normalsize

\section{Conclusion}\label{sec:conclusion}

The architecture has two central claims. First, an AI mathematics system must
navigate problem-sized mathematical worlds rather than treat a large Lean
repository as a bag of files or one prompt. Second, scaling search safely
requires reasoning scope, mutation permission, validator capacity, evidence
class, public-claim permission, and reviewer attention to remain non-fungible
until an explicit fan-in gate. Together these properties make the architecture
a claim-transition operating system with a prover as one component.

The boundaries follow immediately. Questions are not work records. Work
records are not experiments. Experiments are not proofs. Lean proofs are not
interpretations. Interpretations are not public claims. Public claims are not
independent review. Review is not community acceptance.

The private workbench makes the early stages durable and concurrent. Type A
actors mutate live substrate under claims and checks; Type B actors contribute
bounded external reasoning without hidden write authority. Computational
experiments prune and prioritise routes. Lean supplies exact formal authority.
Graphs make a large corpus navigable without collapsing inventory into meaning.
Comparator checks exact statement interfaces. Local review selection directs
scarce review attention, and the external Palomar registry checks and filters
what is submitted to it. The public repository then packages source, claims, papers, and
release receipts into a clone that stands on its own.

The useful result is not an automatic mathematician and not a universal
publication verifier. It is a system in which the evidence for a claim, the
person or program allowed to judge it, and the stronger statement that remains
open are all visible at the point where the claim moves forward. That is what
lets an AI-assisted research process scale without turning accumulated output
into accumulated ambiguity.

\appendix
\small
\section{Reproducibility}\label{app:repro}

\paragraph{Public first contact.}
A fresh public clone can reproduce the control card and structural checks:
\begin{verbatim}
python3 scripts/proof_cockpit.py --format card
python3 scripts/proof_cockpit.py --check
\end{verbatim}
The first reads committed public metadata; the second checks the claim
registry, cold-clone contract, and orientation freshness.  Neither checks a
proof.  Formal authority begins with the pinned Lean build named by the card.

\paragraph{Dated navigation counts.}
At the semantic review's snapshot, taken before the eight-problem
consolidation of 2026-08-02 enlarged the corpus to 153,238 declarations, all
151,761 then-live declarations were inventoried and routed, and all
\SemanticAuthoredTheoremLike\ author-written theorem-like
declarations had an exact node link.  Of those,
\SemanticAuthoredInterpreted\ (\SemanticAuthoredPercent\%) participated in
authored mathematical interpretations: \SemanticDirectEvidence\ as exact
proposition evidence and \SemanticContextual\ as bounded certificate- or
module-family context.  The remaining \SemanticStructuralOnly\ were linked only
through exact source-module and normalised-signature families, not authored
mathematical paraphrases.  Every declaration selected for a public claim had
an authored route.  The command
\texttt{python3 scripts/query\_semantic.py coverage} derives these volatile
navigation counts and checks their references.  They do not measure semantic
review quality or public-claim completeness.

The paper inventory, \path{docs/publication_contract.json}, records source
and PDF cryptographic hashes and validation commands. The worked-example evidence is recorded in
\path{docs/publication_evidence.json}. The reconstruction file
\path{experiments/publication_mutations.json} specifies the deliberately false
edits and their limitations. Comparator's public configuration and receipt are
named by \path{verification/comparator.json}; the local review-selection readiness and
ranking surfaces are named by the corresponding files under \path{docs/}.
These artefacts identify what was checked. They do not interpret unrestricted
prose or confer external acceptance.

The private system is described here at the architectural level because it is
production provenance, not a dependency of the public result. Private paths,
operator material, unreleased work, and private ledgers are neither required
nor granted authority by this paper. The public checkout remains the inspection
and replay boundary.

{\tiny
\begin{multicols}{2}
\begin{thebibliography}{20}
\setlength{\itemsep}{0pt}
\setlength{\parsep}{0pt}
\bibitem{lean4} L.~de Moura and S.~Ullrich, \emph{The Lean 4 Theorem Prover
and Programming Language}, in \emph{Automated Deduction---CADE 28}, Lecture
Notes in Computer Science 12699, 2021, pp.~625--635,
\href{https://doi.org/10.1007/978-3-030-79876-5_37}{DOI}.
\bibitem{erdosgraham} P.~Erd\H{o}s and R.~L.~Graham, \emph{Old and New
Problems and Results in Combinatorial Number Theory}, Monographies de
L'Enseignement Math\'ematique 28, 1980, p.~61.
\bibitem{leanblueprint} P.~Massot, \emph{leanblueprint}, plasTeX plugin for
Lean formalisation blueprints, 2020,
\href{https://github.com/PatrickMassot/leanblueprint}{software repository}.
\bibitem{leanarchitect} T.~Zhu, P.~Monticone, S.~Welleck, and J.~Avigad,
\emph{LeanArchitect: Automating Blueprint Generation for Humans and AI},
in \emph{17th International Conference on Interactive Theorem Proving},
LIPIcs 382, 2026, pp.~25:1--25:16.
\bibitem{leanatlas} B.~Yanahama and A.~Sannai, \emph{Lean Atlas: An
Integrated Proof Environment for Scalable Human--AI Collaborative
Formalization}, 2026, \href{https://doi.org/10.48550/arXiv.2604.16347}{arXiv}.
\bibitem{econcs} N.~Garg, \emph{EconCSLib: AI-Assisted Lean Formalization for
Economics \& Computation Research}, 2026,
\href{https://doi.org/10.48550/arXiv.2606.13306}{arXiv}.
\bibitem{isadof} A.~D.~Brucker and B.~Wolff, \emph{Isabelle/DOF: Design and
Implementation}, in \emph{Software Engineering and Formal Methods}, 2019,
pp.~275--293.
\bibitem{gotel} O.~C.~Z.~Gotel and A.~C.~W.~Finkelstein, \emph{An analysis
of the requirements traceability problem}, Proc. First IEEE International
Conference on Requirements Engineering, 1994, pp.~94--101.
\bibitem{demillo} R.~A.~DeMillo, R.~J.~Lipton, and F.~G.~Sayward, \emph{Hints
on test data selection}, IEEE Computer 11(4), 1978, pp.~34--41.
\bibitem{jiaharman} Y.~Jia and M.~Harman, \emph{An analysis and survey of the
development of mutation testing}, IEEE Transactions on Software Engineering
37(5), 2011, pp.~649--678.
\bibitem{gsn} SCSC Assurance Case Working Group, \emph{Goal Structuring
Notation Community Standard, Version 3}, SCSC-141C, May 2021.
\bibitem{taoai} T.~Tao, \emph{Mathematics in the age of AI}, preprint, 2026,
\href{https://doi.org/10.48550/arXiv.2608.16753}{arXiv}.
\bibitem{mathlib} Lean community, \emph{Contributing to mathlib},
\href{https://leanprover-community.github.io/contribute/index.html}%
{contributor guide}, accessed August 2026.
\bibitem{leandojo} K.~Yang et al., \emph{LeanDojo: Theorem Proving with
Retrieval-Augmented Language Models}, NeurIPS 2023.
\bibitem{pantograph} J.~Storrs et al., \emph{Pantograph: A Machine-to-Machine
Interaction Interface for Advanced Theorem Proving, High Level Reasoning, and
Data Extraction in Lean 4}, 2024,
\href{https://doi.org/10.48550/arXiv.2410.16429}{arXiv}.
\bibitem{openprover} M.~Kripner and M.~Straka, \emph{OpenProver: Agentic and
Interactive Theorem Proving with Lean 4}, 2026,
\href{https://doi.org/10.48550/arXiv.2607.09217}{arXiv}.
\bibitem{agenthunt} C.~E.~Brown, C.~Kaliszyk, and J.~Urban, \emph{Agent Hunt:
Bounty Based Collaborative Autoformalization With LLM Agents}, 2026,
\href{https://doi.org/10.48550/arXiv.2603.06737}{arXiv}.
\bibitem{dreamprover} Y.~Zhang et al., \emph{DreamProver: Evolving
Transferable Lemma Libraries via a Wake--Sleep Theorem-Proving Agent}, 2026,
\href{https://doi.org/10.48550/arXiv.2604.26311}{arXiv}.
\bibitem{hep} I.~Takahara and T.~Mizoguchi, \emph{Toward Auditable AI
Scientists: A Hypothesis Evolution Protocol for LLM Agents}, 2026,
\href{https://doi.org/10.48550/arXiv.2607.09195}{arXiv}.
\bibitem{symposium} D.~Pratt, \emph{Symposium: Trust via Auditable Records for
Communities of AI Scientist Agents}, 2026,
\href{https://doi.org/10.48550/arXiv.2608.19511}{arXiv}.
\bibitem{eurekagent} A.~Xin et al., \emph{EurekAgent: Agent Environment
Engineering is All You Need for Autonomous Scientific Discovery}, 2026,
\href{https://doi.org/10.48550/arXiv.2606.13662}{arXiv}.
\end{thebibliography}
\end{multicols}
}
\end{document}
